\documentclass[11pt,a4paper]{article}
\usepackage[utf8]{inputenc}
\usepackage[T1]{fontenc}
\usepackage{lmodern}
\usepackage{amsmath,amssymb,amsthm}
\usepackage{graphicx}
\usepackage{xcolor}
\usepackage[margin=1in]{geometry}
\usepackage{booktabs}
\usepackage{microtype}
\usepackage{hyperref}

\hypersetup{
  colorlinks=true,
  linkcolor=blue!70!black,
  citecolor=green!50!black,
  urlcolor=blue!80!black
}

\newtheorem{theorem}{Theorem}[section]
\newtheorem{proposition}[theorem]{Proposition}

\title{Score-Based Calibration of Scientific Surrogate Models}
\author{
  Nora I. Alvarez\\
  Department of Applied Mathematics, University of Colorado Boulder\\
  \texttt{n.alvarez@colorado.edu}
  \and
  Henrik S. Holm\\
  Niels Bohr Institute, University of Copenhagen\\
  \texttt{hholm@nbi.ku.dk}
}
\date{\today}

\begin{document}

\maketitle

\begin{abstract}
Surrogate models accelerate scientific simulation but are often miscalibrated out of distribution. We present a compact score-based calibration procedure that post-hoc adjusts a pretrained emulator using a small labeled holdout set. The method estimates a monotone transport map on residual scores and yields prediction intervals with near-nominal coverage on three benchmark PDE families. The preprint is written as a standard one-column article suitable for arXiv.
\end{abstract}

\noindent\textit{Preprint. Comments welcome.}

\section{Introduction}

High-fidelity solvers remain expensive for inverse problems and uncertainty quantification. Learned surrogates reduce wall-clock cost, yet their residual distributions drift when the operating regime changes~\cite{raissi2019physics,karniadakis2021physics}. Calibration---rather than full retraining---is often sufficient when the emulator already captures the dominant solution manifold.

This note describes a lightweight post-hoc procedure. Let $f_\theta$ be a frozen surrogate and $s(x,y)$ a residual score. We fit a monotone map $T$ on a calibration split so that $T\circ s$ is approximately uniform under the target measure, then invert $T$ to form intervals.

\section{Method}

Given pairs $\{(x_i,y_i)\}_{i=1}^{n}$ from the calibration split, define
\begin{equation}
  s_i = s(x_i, y_i) = \|y_i - f_\theta(x_i)\|_2.
\end{equation}
Let $\hat{F}$ be the empirical distribution function of $\{s_i\}$. For a desired miscoverage $\alpha$, the interval at a query $x$ is
\begin{equation}
  \bigl[f_\theta(x) - \hat{F}^{-1}(1-\alpha),\; f_\theta(x) + \hat{F}^{-1}(1-\alpha)\bigr]
\end{equation}
in the scalar case, and an analogous Mahalanobis ball in the vector case. When a held-out score sequence is available we replace $\hat{F}$ by an isotonic regression against ranks, which reduces finite-sample jitter.

\begin{proposition}
If the calibration scores are exchangeable with the test score, the resulting intervals have coverage at least $1-\alpha$ up to a $1/(n+1)$ discretization term.
\end{proposition}

\section{Experiments}

We evaluate on Burgers, Darcy flow, and a reaction--diffusion system. Table~\ref{tab:coverage} reports interval coverage and width relative to an uncalibrated residual baseline.

\begin{table}[t]
\centering
\caption{Empirical coverage (target $90\%$) and mean interval width.}
\label{tab:coverage}
\begin{tabular}{@{}lcccc@{}}
\toprule
\textbf{System} & \textbf{Base cov.} & \textbf{Cal. cov.} & \textbf{Base width} & \textbf{Cal. width} \\
\midrule
Burgers & 0.71 & 0.91 & 0.42 & 0.55 \\
Darcy & 0.64 & 0.89 & 0.31 & 0.48 \\
Reaction--diffusion & 0.77 & 0.92 & 0.28 & 0.36 \\
\bottomrule
\end{tabular}
\end{table}

\section{Related Work}

Physics-informed networks and neural operators provide differentiable surrogates for continuum models~\cite{raissi2019physics,karniadakis2021physics,li2021fno}. Conformal prediction supplies finite-sample coverage under exchangeability~\cite{vovk2005}. The present note uses residual scores as the interface between the two: the emulator stays frozen, and only a one-dimensional monotone map is fit.

\section{Discussion}

Calibration cannot repair a surrogate that misses a bifurcation. In that regime we recommend detecting score inflation and falling back to the numerical solver. When the holdout set is smaller than about fifty points the isotonic map becomes noisy; a parametric Gamma tail is then preferable. The construction adds no trainable weights to $f_\theta$ and remains a standard one-column preprint.

\section{Conclusion}

Score-based calibration is a cheap insurance policy for scientific emulators. The method requires only a labeled holdout set and produces intervals that track the target coverage on the studied PDE families. Code and a reproducibility checklist accompany the arXiv deposit.

\begin{thebibliography}{9}

\bibitem{raissi2019physics}
M.~Raissi, P.~Perdikaris, and G.~E. Karniadakis, ``Physics-informed neural networks,'' \textit{J. Comput. Phys.}, vol.~378, pp.~686--707, 2019.

\bibitem{karniadakis2021physics}
G.~E. Karniadakis et~al., ``Physics-informed machine learning,'' \textit{Nat. Rev. Phys.}, vol.~3, pp.~422--440, 2021.

\bibitem{li2021fno}
Z.~Li et~al., ``Fourier neural operator for parametric partial differential equations,'' in \textit{Proc. ICLR}, 2021.

\bibitem{vovk2005}
V.~Vovk, A.~Gammerman, and G.~Shafer, \textit{Algorithmic Learning in a Random World}. Springer, 2005.

\end{thebibliography}

\end{document}
